\documentclass[a4paper,14pt,oneside]{extarticle}

\usepackage[utf8]{inputenc}
\usepackage[T2A]{fontenc}
\usepackage[english,russian]{babel}
\usepackage[left=3cm,right=1.5cm,top=2cm,bottom=2cm]{geometry}
\usepackage{setspace}
\usepackage{amsmath}
\usepackage{graphicx}
\usepackage{float}
\usepackage{array}
\usepackage{longtable}
\usepackage{verbatim}
\usepackage[hidelinks]{hyperref}

\setstretch{1.35}
\setlength{\parindent}{1.25cm}
\setlength{\parskip}{0pt}
\addto\captionsrussian{
  \renewcommand{\contentsname}{Содержание}
  \renewcommand{\figurename}{Рисунок}
}

\begin{document}
\selectlanguage{russian}

\begin{titlepage}
\begin{center}
\textbf{МИНИСТЕРСТВО НАУКИ И ВЫСШЕГО ОБРАЗОВАНИЯ РФ}\\
\vspace{0.4cm}
\textbf{[Наименование образовательной организации]}\\
\textbf{[Наименование кафедры]}

\vfill

\textbf{\Large ОТЧЁТ}\\
\vspace{0.4cm}
\textbf{\large по лабораторной работе № 1}\\
\vspace{0.4cm}
\textbf{\large «Генерация паролей»}

\vfill

\begin{flushright}
\begin{minipage}{0.48\textwidth}
Выполнил: Сечейко Н.В.\\
\vspace{0.3cm}
Проверил(а): \underline{\hspace{4cm}}
\end{minipage}
\end{flushright}

\vfill

\textbf{Минск}\\
\textbf{2026}
\end{center}
\end{titlepage}

\tableofcontents
\newpage

\section{Цель и содержание работы}

Цель работы --- разработать программу генерации паролей, исследовать
равномерность распределения символов выбранного алфавита, оценить среднее
время полного перебора пароля и сформулировать рекомендации по выбору длины
пароля.

В соответствии с методическими указаниями программа должна:
\begin{enumerate}
\item генерировать строку заданной пользователем длины из алфавита варианта;
\item проверять равномерность распределения символов с помощью частотной
гистограммы;
\item вычислять среднее время подбора пароля, выбранного из сгенерированной
строки;
\item строить график зависимости среднего времени подбора от длины пароля;
\item формулировать практические рекомендации по выбору пароля.
\end{enumerate}

\section{Теоретические сведения}

\subsection{Идентификация и аутентификация}

Идентификация представляет собой установление идентификатора субъекта в
информационной системе. Аутентификация проверяет принадлежность субъекту
предъявленного идентификатора и подтверждает его подлинность.

Одним из наиболее распространённых способов аутентификации является
использование паролей. Надёжность случайного пароля зависит от мощности
алфавита и его длины. Если пароль составляется из символов алфавита мощности
$A$ и имеет длину $L$, количество возможных паролей равно

\[
N_{\text{паролей}} = A^L.
\]

При полном переборе пароль в среднем находится после проверки половины
вариантов:

\[
N_{\text{ср}} = \frac{A^L+1}{2}.
\]

Если средство атаки проверяет $V$ паролей в секунду, среднее время подбора
составляет

\[
T_{\text{ср}} = \frac{A^L+1}{2V}.
\]

Каждый дополнительный символ увеличивает количество вариантов в $A$ раз.
Поэтому увеличение длины случайного пароля является эффективным способом
повышения его стойкости.

\subsection{Равномерность распределения}

Для равномерного генератора вероятность появления каждого символа алфавита
одинакова:

\[
P(x_i)=\frac{1}{A}.
\]

В варианте 26 используется алфавит из 36 символов, поэтому ожидаемая частота
каждого символа равна

\[
\frac{1}{36}\approx 0{,}02778=2{,}778\%.
\]

Для строки длиной $N$ ожидаемое количество вхождений одного символа:

\[
n_i=\frac{N}{36}.
\]

Из-за случайного характера генерации фактические частоты не обязаны совпадать
с ожидаемыми точно. Для проверки строится гистограмма частот.

\section{Исходные данные и вариант}

Используется вариант 26. Согласно таблице вариантов ему соответствует
алфавит из строчных латинских букв и арабских цифр:

\begin{center}
\texttt{abcdefghijklmnopqrstuvwxyz0123456789}
\end{center}

Мощность алфавита:

\[
A=26+10=36.
\]

\section{Описание разработанной программы}

Программа написана на языке C++. Для выполнения требований методички
используются функции \texttt{rand()} и \texttt{srand()}. Начальное значение
генератора инициализируется текущим временем.

Пользователь вводит:
\begin{itemize}
\item $N$ --- длину генерируемой строки;
\item $L$ --- длину пароля, выбираемого из этой строки.
\end{itemize}

После генерации программа выводит сокращённое представление строки и
выбранный пароль. Для каждого из 36 символов подсчитывается количество
вхождений и строится текстовая гистограмма.

Для оценки времени полного перебора в программе принято учебное значение
$V=10^9$ попыток в секунду. Время рассчитывается по формуле

\[
T_{\text{ср}}=\frac{36^L+1}{2\cdot 10^9}.
\]

График зависимости времени от длины пароля выводится в консоли в
логарифмическом масштабе. Логарифмический масштаб необходим, поскольку
значения времени изменяются от наносекунд до миллионов лет.

\section{Ход выполнения и результаты}

Для экспериментальной проверки были выбраны:

\[
N=1000,\qquad L=16.
\]

Ожидаемое количество каждого символа:

\[
\frac{1000}{36}\approx27{,}78.
\]

Пример сгенерированного пароля:

\begin{center}
\texttt{674wlovx0jk5s9t7}
\end{center}

Фактическая строка и частотная гистограмма были получены при запуске
разработанной программы.

\subsection{Расчёт для пароля длиной 16}

Количество возможных паролей:

\[
36^{16}\approx7{,}959\cdot10^{24}.
\]

Среднее количество попыток:

\[
\frac{36^{16}+1}{2}\approx3{,}979\cdot10^{24}.
\]

При скорости $10^9$ попыток в секунду:

\[
T_{\text{ср}}\approx1{,}26\cdot10^8\text{ лет}.
\]

Полученный результат показывает экспоненциальный рост стойкости при
увеличении длины пароля.

\subsection{Зависимость от длины пароля}

\begin{center}
\begin{tabular}{|c|c|}
\hline
\textbf{Длина $L$} & \textbf{Среднее время при $V=10^9$ попыток/с}\\
\hline
6 & 1,09 с\\
\hline
8 & 23,51 мин\\
\hline
10 & 21,16 суток\\
\hline
12 & 75,08 года\\
\hline
14 & 97297 лет\\
\hline
16 & $1{,}26\cdot10^8$ лет\\
\hline
\end{tabular}
\end{center}

Текстовый график, построенный программой, показывает монотонный
экспоненциальный рост среднего времени подбора. При добавлении одного
символа время увеличивается примерно в 36 раз.

\section{Практические рекомендации}

Рекомендации зависят от ценности информации, скорости атакующего и времени,
в течение которого информация должна оставаться защищённой.

\begin{itemize}
\item Для малозначимой учётной записи при скорости $10^9$ попыток/с следует
использовать случайный пароль длиной не менее 10 символов.
\item Для почты и личных данных рекомендуется длина не менее 14 символов с
учётом возможности мощного офлайн-перебора.
\item Для финансовой информации и коммерческой тайны рекомендуется длина не
менее 15--16 случайных символов.
\item Нельзя использовать имена, даты рождения и словарные слова, поскольку
для них применяются словарные атаки.
\item Для реальных систем необходимо ограничивать число попыток входа и
хранить не сам пароль, а его защищённое хеш-значение.
\end{itemize}

Использованный в работе \texttt{rand()} предназначен для выполнения учебного
задания. Для реальных паролей следует применять криптографически стойкий
генератор случайных чисел.

\section{Заключение}

В ходе лабораторной работы разработана программа генерации паролей для
варианта 26. Использован алфавит из строчных латинских букв и арабских цифр
мощностью 36 символов. Выполнена проверка частотного распределения символов
с помощью гистограммы, рассчитано среднее время полного перебора и построен
график его зависимости от длины пароля.

Эксперимент показал, что увеличение длины случайного пароля значительно
повышает его стойкость: каждый дополнительный символ увеличивает пространство
перебора в 36 раз. Для защиты важной информации рекомендуется использовать
длинные случайные пароли и дополнять парольную аутентификацию ограничением
числа попыток и безопасным хранением хешей.

\appendix
\section{Исходный код программы}

\verbatiminput{lab1.cpp}

\end{document}
